% ---------------------------------------------------------------------------
% INTRODUCTION --- ICRA draft, style-calibrated to DextrAH-G/RGB and Kang et al.
% Drop-in replacement for \section{INTRODUCTION} in the ieeeconf template.
% Uses the cite keys already in your draft (zhu2025deformable, blanco2024t,
% yin2021modeling, arriola2020modeling, kang2026learning) plus the verified ones
% in docs/paper/refs.bib (pan2020stressmin, vtam, lighttact, curseofprecision,
% wang2025damageless, wang2024softgripper).
%
% Register notes (what makes this ICRA rather than NeurIPS):
%   - short declarative sentences, one claim each; no long subordinate clauses
%   - "However," pivot after the opening challenge sentence
%   - enumerated (i)(ii)(iii) factors and challenges, kept parallel
%   - headline numbers in the Introduction, not deferred to Experiments
%   - contributions as an itemize, each one line, each mapping to a section
%
% PLACEHOLDER: \method --- name the system before submission. A named system
% reads as a contribution; "our method" reads as an experiment.
% ---------------------------------------------------------------------------
\newcommand{\method}{\textsc{Method}}   % TODO: name it

\section{INTRODUCTION}

Manipulating delicate food is a core capability for automation in food processing,
agriculture, and assistive settings, and it remains difficult for current
robots~\cite{wang2025damageless}. Produce tears, cracks, and bruises under modest
compression or tension, and the damage is
irreversible~\cite{zhu2025deformable,blanco2024t}. Handling it well requires getting two
quantities right at once: \emph{where} and in what orientation to grasp the item, and
\emph{how far} to close the gripper. A grasp that straddles a stem or closes on a thin edge
fails at a width that would have been correct elsewhere on the same item. Both quantities
depend on the item's pose, size, shape, and stiffness. From a single depth view the first
three are only partially observable and the fourth is not observable at all: two visually
identical berries at different ripeness demand different grips. The difficulty stems from
three factors: (i) contact and deformation dynamics are hard to model, so the internal stress
that causes damage cannot be assessed at run
time~\cite{yin2021modeling,arriola2020modeling}; (ii) the safety constraint is tight, since
stress must stay below an item-dependent damage threshold; and (iii) shape, size, and
physical properties vary enormously both within a food category and across categories.

Existing approaches address the first two factors and leave the third open. Compliant and
pneumatic end-effectors bound contact pressure mechanically~\cite{wang2024softgripper}, but
compliance is a fixed property of the hardware and does not decide where to grasp. Tactile
and force-controlled grippers close the loop on contact at run time and handle delicate items
well~\cite{kang2026learning,vtam}. However, these policies are trained by imitation from human
demonstrations collected on real hardware, so the cost of acquiring data is unchanged.
Tactile signals are also not reliably simulatable, which closes the escape route to simulated
data, and deformation-based sensors require measurable indentation to produce a signal at
all~\cite{lighttact}. Stress-aware grasp planning targets the damaging quantity
directly~\cite{pan2020stressmin}, but assumes point contacts and scores grasps by resistible
external wrench. We implemented that metric and found it uninformative for our setting: two
point contacts on a smooth rounded object cannot span the moment axes of wrench space, so the
metric returns essentially zero for every candidate grasp, while a cube clears it.

None of these lines supplies coverage of factor (iii), and demonstrations cannot supply it
either. A demonstration is recorded against one item in one pose and cannot be transported to
another, so coverage cost scales with the product of the varying dimensions and multiplies
again across categories. For precision-limited tasks the number of demonstrations required
grows super-exponentially as the tolerance tightens, $\log N \propto 1/(P-c)$
\cite{curseofprecision}. We observe the consequence directly. A policy trained on 50 real
demonstrations does not degrade uniformly on hardware; it fails in a structured way,
misgrasping at particular workspace locations and mis-squeezing at particular object sizes
(Fig.~\ref{fig:coverage}). The bottleneck is coverage of the physical variation, not the
number of trajectories.

We present \method, a sim-to-real pipeline that obtains that coverage from simulation rather
than from human effort. Our key observation is that the quantity which determines gentleness
is unavailable on hardware but exact in a deformable simulator: internal stress. We therefore
use it as a \emph{planning objective} rather than a sensed signal. A stress-aware planner
scores a candidate grasp by the internal stress it induces under a width-controlled,
distributed-pad contact model, matching a real parallel jaw that commands a displacement and
obtains a force as an outcome, and searches the $7$-DoF grasp with CMA-ES subject to holding
the item against gravity. Because the displacement is prescribed, stress is linear in Young's
modulus, so a single finite-element solve per grasp pose covers an entire distribution of
stiffnesses. Randomising material therefore costs almost nothing, which is what makes the
planner a coverage instrument rather than a grasp ranker. We run it over randomised object
pose, size, shape, and material to generate thousands of gentle demonstrations without a human
in the loop, distil them into a point-cloud diffusion policy acting from a single depth view,
and co-train with a small slice of real demonstrations that calibrates the residual sensor
gap. \textbf{[TODO: headline real-robot number + \# objects, stated here as DextrAH does.]}

\noindent Our contributions are:
\begin{itemize}
  \item A stress-aware grasp planner built on a width-controlled, distributed-pad contact
        model, which discriminates among parallel-jaw grasps on organic objects where the
        point-contact force-closure formulation is degenerate.
  \item A stiffness-linearity result that makes randomised material coverage tractable: one
        finite-element solve per pose covers an entire stiffness distribution.
  \item An autonomous demonstration generator that produces gentle, successful grasps across
        randomised object pose, size, shape, and material with no human demonstration.
  \item Real-robot validation across object poses, sizes, shapes, and categories from a single
        depth view, with no tactile or force sensing in the loop.
        \textbf{[TODO: multi-category run is blocking --- see PAPER\_TODO A7.]}
\end{itemize}
